%%%%%%%%%%%%%%%%%%%%%%%%%%%%%%%%%%%%%%%%%
% University Assignment Title Page 
% LaTeX Template
% Version 1.0 (27/12/12)
%
% This template has been downloaded from:
% http://www.LaTeXTemplates.com
%
% Original author:
% WikiBooks (http://en.wikibooks.org/wiki/LaTeX/Title_Creation)
%
% License:
% CC BY-NC-SA 3.0 (http://creativecommons.org/licenses/by-nc-sa/3.0/)
% 
% Instructions for using this template:
% This title page is capable of being compiled as is. This is not useful for 
% including it in another document. To do this, you have two options: 
%
% 1) Copy/paste everything between \begin{document} and \end{document} 
% starting at \begin{titlepage} and paste this into another LaTeX file where you 
% want your title page.
% OR
% 2) Remove everything outside the \begin{titlepage} and \end{titlepage} and 
% move this file to the same directory as the LaTeX file you wish to add it to. 
% Then add \input{./title_page_1.tex} to your LaTeX file where you want your
% title page.
%
%%%%%%%%%%%%%%%%%%%%%%%%%%%%%%%%%%%%%%%%%
%\title{Title page with logo}
%----------------------------------------------------------------------------------------
%	PACKAGES AND OTHER DOCUMENT CONFIGURATIONS
%----------------------------------------------------------------------------------------

\documentclass[12pt]{article}
\usepackage[utf8]{inputenc}
\usepackage[brazil]{babel}
\usepackage{amsmath}
\usepackage{amssymb}
\usepackage{graphicx}
\usepackage[colorinlistoftodos]{todonotes}
\usepackage{enumitem}
\usepackage{subcaption}
\usepackage{geometry}
    \geometry{
    a4paper,
    total={170mm,257mm},
    left=20mm,
    top=20mm,
    }
\usepackage[document]{ragged2e}
\usepackage[toc,page]{appendix}
\usepackage[framed,numbered,autolinebreaks,useliterate]{mcode}
\usepackage{float}

\newcommand{\HRule}{\rule{\linewidth}{0.5mm}} % Defines a new command for the horizontal lines, change thickness here
\renewcommand\appendixpagename{Apêndice: Identidades trigonométricas}

\begin{document}
%----------------------------------------------------------------------------------------
%	HEADING SECTIONS
%----------------------------------------------------------------------------------------

\begin{minipage}{0.15\textwidth}
    \begin{flushleft}
        \includegraphics[width=\linewidth]{logo.png}
    \end{flushleft}
\end{minipage}
~
\begin{minipage}{0.75\textwidth}
    \textsc{Universidade de São Paulo}\\
    \textsc{Escola de Engenharia de São Carlos}\\
    \textsc{Departamento de Engenharia Elétrica e de Computação}\\
    \textsc{\textbf{SEL0343 | Processamento Digital de Sinais}}\\
    \textsc{\textbf{Prof. Dr. Emiliano R. Martins}}
\end{minipage}

\begin{flushright}
    {\small \today}
\end{flushright}

%----------------------------------------------------------------------------------------
%	AUTHOR SECTION
%----------------------------------------------------------------------------------------

\begin{minipage}{0.52\textwidth}
\begin{flushleft} \large
\emph{Aluno:}\\
\textbf{Luís Filipe Vasconcelos Peres} % Your name
\end{flushleft}
\end{minipage}
~
\begin{minipage}{0.45\textwidth}
\begin{flushright} \large
\emph{Nº USP:} \\
\textbf{10310641} % Supervisor's Name
\end{flushright}
\end{minipage}\\

%----------------------------------------------------------------------------------------
%	TITLE SECTION
%----------------------------------------------------------------------------------------
\center
\HRule \\[0.4cm]
{ \LARGE \bfseries Prova única}\\ % Title of your document
\HRule \\[0.2cm]

%----------------------------------------------------------------------------------------
%	CONTENT SECTION
%----------------------------------------------------------------------------------------
\justify
\section{}

    
\appendix
\appendixpage

\begin{equation}
    \begin{split}
        \sin(ax + ay) & = \sin(ax) \cdot \cos(ay) + \sin(ay) \cdot \cos(ax)\\
        \sin(ax - ay) & = \sin(ax) \cdot \cos(ay) - \sin(ay) \cdot \cos(ax)\\
        \sin(ax + ay) + \sin(ax - ay) & = 2 \cdot \sin(ax) \cdot \cos(ay)    
    \end{split}
\end{equation}

\begin{equation}\label{eq:2}
    \begin{split}
        \cos( \alpha - \beta ) + \cos( \alpha + \beta ) = &~(\cos\alpha \cdot \cos\beta - \sin\alpha \cdot \sin\beta) ~+ \\
        &~ (\cos\alpha \cdot \cos\beta + \sin\alpha \cdot \sin\beta)\\
        = & ~2 \cdot \cos\alpha \cdot \cos\beta
    \end{split}
\end{equation}

\begin{equation}\label{eq:3}
    \begin{split}
        \cos(2\alpha) &= \cos^2\alpha - \sin^2\alpha\\
        &= \cos^2\alpha - (1 - \cos^2\alpha)\\
        &= 2\cdot \cos^2\alpha - 1\\
        \Rightarrow ~2\cdot \cos^2\alpha &= 1 + \cos(2\alpha)\\
        \therefore ~\cos^2\alpha &= \frac{1 + \cos(2\alpha)}{2}
    \end{split}
\end{equation}

\end{document}